% 제4장: 실험 결과 및 분석 / 4.1 초안
% article 계열의 제4절 아래에 삽입하는 본문 조각이다.
% 근거: results/koblex-reproduction-v1/protocol-lock.json,
%       results/koblex-reproduction-v1/chapter4-prep-v1/verification-audit.json,
%       results/koblex-reproduction-v1/chapter4-prep-v1/A1/numeric-ledger.json

\subsection{실험 조건 및 비교 기준}
\label{sec:ch4_conditions}

본 장에서는 KoBLEX의 공식 한국어 226문항을 대상으로, ParSeR의
CE(cross-encoder)를 Jev로 교체했을 때 최종 근거 선택과 답변 품질에
나타나는 차이를 분석한다. 비교 조건을 표~\ref{tab:ch4_conditions}에 정리하였다.
공개 코드의 후보 목록 참조 인덱스 문제를 수정한 CE를 주 기준선으로
두고, 동일 파이프라인에서 점수화만 Jev로 교체한 조건과 대응 비교한다.
인덱스 문제를 보존한 공개 코드 CE는 재현 점검을 위한 참조 조건으로 제시한다.

\begin{table}[htbp]
    \centering
    \caption{비교 조건의 구성}
    \label{tab:ch4_conditions}
    \small
    \setlength{\tabcolsep}{4pt}
    \begin{tabular}{@{}llll@{}}
        \hline
        조건 & 점수화 & 선택 인덱스 & 역할 \\
        \hline
        공개 코드 CE & BGE & 문제 보존 & 재현 참조 \\
        수정 CE      & BGE & 수정      & 주 기준선 \\
        Jev 교체     & Jev & 수정      & 주 비교 \\
        \hline
    \end{tabular}
\end{table}

세 조건은 동일한 예상 법령과 BM25 후보를 공유한다. 예상 법령별 후보
100개를 재정렬한 뒤 상위 10개를 공통 LLM 선택기에 제시하고,
각 목록에서 한 개를 선택하여 답변을 생성하였다. 최종 근거 수는 생성된
예상 법령 수에 따라 달라지며, 정답 hop 정보는 선택에 사용하지 않았다.
선택ㆍ생성 프롬프트와 평가 기준도 동일하게 유지하였다.
CE는 \texttt{bge-reranker-v2-m3-ko}, Jev는 \texttt{jev-1.14.0}을
사용하였으며, 선택ㆍ답변 생성과 LF-Eval의 평가 모델은
\texttt{gpt-4o-2024-08-06}으로 통일하였다.

정답 근거의 일부가 corpus와 정확히 매핑되지 않는 29문항도 제외하지
않고 전체 226문항을 주 분석의 분모로 유지하였다. 이하의 비교는 기존
시험 결과를 확인한 이후 수행한 탐색적 재현 분석으로 보고한다.
